\documentclass[UTF8,12pt]{ctexart}
\usepackage[a4paper,margin=24mm]{geometry}
\usepackage{amsmath,amssymb,bm,graphicx,adjustbox}
\newcommand{\fitmath}[1]{\begin{adjustbox}{max width=\linewidth}$\displaystyle #1$\end{adjustbox}}
\setlength{\parindent}{0pt}
\setlength{\parskip}{.7em}
\sloppy
\allowdisplaybreaks
\begin{document}
\section*{1997 年考研数学三 · 真题}
\subsection*{1997 年全国硕士研究生招生考试试题}

\subsection*{一、填空题（本题共 5 小题，每小题 3 分，满分 15 分）}

（1）设 \(\displaystyle y=f(\ln x)e^{f(x)}\)，其中 \(\displaystyle f\) 可微，则 \(\displaystyle dy=\)\_\_\_\_\_\_。


（2）若 \(\displaystyle f(x)=\dfrac{\displaystyle 1}{\displaystyle 1+x^2}+\sqrt{1-x^2}\displaystyle \int_0^1f(x)\,dx\)，则 \(\displaystyle \int_0^1f(x)\,dx=\)\_\_\_\_\_\_。


（3）差分方程 \(\displaystyle y_{t+1}-y_t=t2^t\) 的通解为\_\_\_\_\_\_。


（4）若二次型 \(\displaystyle f(x_1,\allowbreak x_2,\allowbreak x_3)=2x_1^2+x_2^2+x_3^2+2x_1x_2+x_1x_3+tx_2x_3\) 是正定的，则 \(\displaystyle t\) 的取值范围是\_\_\_\_\_\_。


（5）设随机变量 \(\displaystyle X\) 和 \(\displaystyle Y\) 相互独立，且都服从正态分布 \(\displaystyle N(0,\allowbreak 3^2)\)，而 \(\displaystyle X_1,\allowbreak \ldots,\allowbreak X_9\) 和 \(\displaystyle Y_1,\allowbreak \ldots,\allowbreak Y_9\) 分别是来自总体 \(\displaystyle X\) 和 \(\displaystyle Y\) 的简单随机样本，则统计量 \(\displaystyle U=\dfrac{\displaystyle X_1+\cdots+X_9}{\displaystyle \sqrt{Y_1^2+\cdots+Y_9^2}}\) 服从\_\_\_\_\_\_分布（2 分），参数为\_\_\_\_\_\_（1 分）。


\subsection*{二、选择题（本题共 5 小题，每小题 3 分，满分 15 分）}

（1）设函数 \(\displaystyle f(x)=\displaystyle \int_0^{1-\cos x}\sin(t^2)\,dt,\allowbreak \ g(x)=\dfrac{\displaystyle x^5}{\displaystyle 5}+\dfrac{\displaystyle x^6}{\displaystyle 6}\)，则当 \(\displaystyle x\to0\) 时，\(\displaystyle f(x)\) 是 \(\displaystyle g(x)\) 的（　）。


（A）低阶无穷小。 （B）高阶无穷小。 （C）等价无穷小。 （D）同阶但不等价的无穷小。


（2）若函数 \(\displaystyle f(-x)=f(x)\)（\(\displaystyle -\infty<x<+\infty\)），在 \(\displaystyle (-\infty,\allowbreak 0)\) 内 \(\displaystyle f^{\prime}(x)>0\)，且 \(\displaystyle f^{\prime\prime}(x)<0\)，则在 \(\displaystyle (0,\allowbreak +\infty)\) 内有（　）。


（A）\(\displaystyle f^{\prime}(x)>0,\allowbreak f^{\prime\prime}(x)<0\)。 （B）\(\displaystyle f^{\prime}(x)>0,\allowbreak f^{\prime\prime}(x)>0\)。 （C）\(\displaystyle f^{\prime}(x)<0,\allowbreak f^{\prime\prime}(x)<0\)。 （D）\(\displaystyle f^{\prime}(x)<0,\allowbreak f^{\prime\prime}(x)>0\)。


（3）设向量组 \(\displaystyle \alpha_1,\allowbreak \alpha_2,\allowbreak \alpha_3\) 线性无关，则下列向量组中，线性无关的是（　）。


（A）\(\displaystyle \alpha_1+\alpha_2,\allowbreak \alpha_2+\alpha_3,\allowbreak \alpha_3-\alpha_1\)。 （B）\(\displaystyle \alpha_1+\alpha_2,\allowbreak \alpha_2+\alpha_3,\allowbreak \alpha_1+2\alpha_2+\alpha_3\)。


（C）\(\displaystyle \alpha_1+2\alpha_2,\allowbreak 2\alpha_2+3\alpha_3,\allowbreak 3\alpha_3+\alpha_1\)。 （D）\(\displaystyle \alpha_1+\alpha_2+\alpha_3,\allowbreak 2\alpha_1-3\alpha_2+22\alpha_3,\allowbreak 3\alpha_1+5\alpha_2-5\alpha_3\)。


（4）设 \(\displaystyle A,\allowbreak B\) 为同阶可逆矩阵，则（　）。


（A）\(\displaystyle AB=BA\)。 （B）存在可逆矩阵 \(\displaystyle P\)，使 \(\displaystyle P^{-1}AP=B\)。 （C）存在可逆矩阵 \(\displaystyle C\)，使 \(\displaystyle C^{\mathrm T}AC=B\)。 （D）存在可逆矩阵 \(\displaystyle P\) 和 \(\displaystyle Q\)，使 \(\displaystyle PAQ=B\)。


（5）设两个随机变量 \(\displaystyle X\) 与 \(\displaystyle Y\) 相互独立，且同分布；\(\displaystyle P\{X=-1\}=P\{Y=-1\}=\dfrac{\displaystyle 1}{\displaystyle 2},\allowbreak \ P\{X=1\}=P\{Y=1\}=\dfrac{\displaystyle 1}{\displaystyle 2}\)，则下列各式中成立的是（　）。


（A）\(\displaystyle P\{X=Y\}=\dfrac{\displaystyle 1}{\displaystyle 2}\)。 （B）\(\displaystyle P\{X=Y\}=1\)。


\subsection*{二、选择题（续）}

（5）（续）（C）\(\displaystyle P\{X+Y=0\}=\dfrac{\displaystyle 1}{\displaystyle 4}\)。 （D）\(\displaystyle P\{XY=1\}=\dfrac{\displaystyle 1}{\displaystyle 4}\)。


三、（本题满分 6 分）在经济学中，称函数 \(\displaystyle Q(x)=A[\delta K^{-x}+(1-\delta)L^{-x}]^{-1/x}\) 为固定替代弹性生产函数，而称函数 \(\displaystyle \bar Q=AK^\delta L^{1-\delta}\) 为 Cobb-Douglas 生产函数（简称 C-D 生产函数）。试证明：当 \(\displaystyle x\to0\) 时，固定替代弹性生产函数变为 C-D 生产函数，即 \(\displaystyle \lim_{x\to0}Q(x)=\bar Q\)。


四、（本题满分 5 分）设 \(\displaystyle u=f(x,\allowbreak y,\allowbreak z)\) 有连续偏导数，\(\displaystyle y=y(x)\) 和 \(\displaystyle z=z(x)\) 分别由方程 \(\displaystyle e^{xy}-y=0\) 和 \(\displaystyle e^z-xz=0\) 所确定，求 \(\displaystyle \dfrac{\displaystyle du}{\displaystyle dx}\)。


五、（本题满分 6 分）一商家销售某种商品的价格满足关系 \(\displaystyle p=7-0.2x\)（万元／吨），\(\displaystyle x\) 为销售量（单位：吨），商品的成本函数是 \(\displaystyle C=3x+1\)（万元）。（1）若每销售一吨商品，政府要征税 \(\displaystyle t\)（万元），求该商家获取最大利润时的销售量；（2）\(\displaystyle t\) 为何值时，政府税收总额最大？


六、（本题满分 6 分）设函数 \(\displaystyle f(x)\) 在 \(\displaystyle [0,\allowbreak +\infty)\) 上连续、单调不减且 \(\displaystyle f(0)\geq0\)。试证函数 \(\displaystyle F(x)=\begin{cases}\dfrac{\displaystyle 1}{\displaystyle x}\displaystyle \int_0^x t^nf(t)\,dt,\allowbreak &x>0,\allowbreak \\0,\allowbreak &x=0\end{cases}\) 在 \(\displaystyle [0,\allowbreak +\infty)\) 上连续且单调不减（其中 \(\displaystyle n>0\)）。


七、（本题满分 6 分）从点 \(\displaystyle P_0(1,\allowbreak 0)\) 作 \(\displaystyle x\) 轴的垂线，交抛物线 \(\displaystyle y=x^2\) 于点 \(\displaystyle Q_0(1,\allowbreak 1)\)；再从 \(\displaystyle Q_0\) 作这条抛物线的切线与 \(\displaystyle x\) 轴交于 \(\displaystyle P_1\)，然后从 \(\displaystyle P_1\) 作 \(\displaystyle x\) 轴的垂线，交抛物线于点 \(\displaystyle Q_1\)，依次重复上述过程得到一系列的点 \(\displaystyle P_1,\allowbreak Q_1;P_2,\allowbreak Q_2;\ldots;P_n,\allowbreak Q_n;\ldots\)。（1）求 \(\displaystyle \overline{OP_n}\)；（2）求级数 \(\displaystyle \overline{Q_1P_1}+\overline{Q_2P_2}+\cdots+\overline{Q_nP_n}+\cdots\) 的和，其中 \(\displaystyle n\)（\(\displaystyle n\geq1\)）为自然数，而 \(\displaystyle \overline{M_1M_2}\) 表示点 \(\displaystyle M_1\) 与 \(\displaystyle M_2\) 之间的距离。


八、（本题满分 6 分）设函数 \(\displaystyle f(t)\) 在 \(\displaystyle [0,\allowbreak +\infty)\) 上连续，且满足方程 \(\displaystyle f(t)=e^{4\pi t^2}+\displaystyle \iint_{x^2+y^2\leq4t^2}f\left(\dfrac{\displaystyle 1}{\displaystyle 2}\sqrt{x^2+y^2}\right)\,dx\,dy\)，求 \(\displaystyle f(t)\)。


九、（本题满分 6 分）设 \(\displaystyle A\) 为 \(\displaystyle n\) 阶非奇异矩阵，\(\displaystyle \alpha\) 为 \(\displaystyle n\) 维列向量，\(\displaystyle b\) 为常数。记分块矩阵 \(\displaystyle P=\begin{pmatrix}E&O\\-\alpha^{\mathrm T}A^*&|A|\end{pmatrix},\allowbreak \ Q=\begin{pmatrix}A&\alpha\\\alpha^{\mathrm T}&b\end{pmatrix}\)，其中 \(\displaystyle A^*\) 是矩阵 \(\displaystyle A\) 的伴随矩阵，\(\displaystyle E\) 为 \(\displaystyle n\) 阶单位矩阵。（1）计算并化简 \(\displaystyle PQ\)；（2）证明：矩阵 \(\displaystyle Q\) 可逆的充分必要条件是 \(\displaystyle \alpha^{\mathrm T}A^{-1}\alpha\ne b\)。


十、（本题满分 10 分）设 3 阶实对称矩阵 \(\displaystyle A\) 的特征值是 1、2、3；矩阵 \(\displaystyle A\) 的属于特征值 1、2 的特征向量分别是 \(\displaystyle \alpha_1=(-1,\allowbreak -1,\allowbreak 1)^{\mathrm T},\allowbreak \ \alpha_2=(1,\allowbreak -2,\allowbreak -1)^{\mathrm T}\)。（1）求 \(\displaystyle A\) 的属于特征值 3 的特征向量；（2）求矩阵 \(\displaystyle A\)。


十一、（本题满分 7 分）假设随机变量 \(\displaystyle X\) 的绝对值不大于 1，\(\displaystyle P\{X=-1\}=\dfrac{\displaystyle 1}{\displaystyle 8},\allowbreak \ P\{X=1\}=\dfrac{\displaystyle 1}{\displaystyle 4}\)，在事件 \(\displaystyle \{-1<X<1\}\) 出现的条件下，\(\displaystyle X\) 在 \(\displaystyle (-1,\allowbreak 1)\) 内的任一子区间上取值的条件概率与该子区间长度成正比。试求 \(\displaystyle X\) 的分布函数。


十二、（本题满分 6 分）游客乘电梯从底层到电视塔顶层观光；电梯于每个整点的第 5 分钟、25 分钟和 55 分钟从底层起运行。假设一游客在早晨八点的第 \(\displaystyle X\) 分钟到达底层候梯处，且 \(\displaystyle X\) 在 \(\displaystyle [0,\allowbreak 60]\) 上均匀分布，求该游客等候时间的数学期望。


十三、（本题满分 6 分）两台同样的自动记录仪，每台无故障工作的时间服从参数为 5 的指数分布。首先开动其中一台，当其发生故障时停用而另一台自行开动。试求两台记录仪无故障工作的总时间 \(\displaystyle T\) 的概率密度 \(\displaystyle f(t)\)、数学期望和方差。

\end{document}
